\documentclass[11pt]{article}
\usepackage[margin=27mm]{geometry}
\usepackage{amsmath,amssymb,amsthm,booktabs,hyperref}
\newtheorem{lemma}{Lemma}\newtheorem{theorem}[lemma]{Theorem}
\newcommand{\PSL}{\operatorname{PSL}}\newcommand{\PGL}{\operatorname{PGL}}\newcommand{\N}{\operatorname{N}}
\title{Exceptional parabolics from their common normalizers}
\author{Human Compression Project}\date{}
\begin{document}\maketitle
\begin{abstract}
A fixed small intersection has at most two exceptional overgroups
of the required type. Counting those overgroups and then forgetting
their order recovers the exceptional-family rows of the Part II
classification. Dihedral second parabolics are uniquely specified by
their rotations; mixed $A_5,S_4$ pairs are specified by an $A_4$
normalizer orbit. The same arguments cover both conjugacy classes
and their fusion, without a numerical group census.
\end{abstract}

\section{Scope and the two-class argument}
Let $G=\PSL(2,q)$ with $q$ odd. For $M\cong S_4$ assume
$q=p$ is prime and $q\equiv\pm1\pmod8$. For $M\cong A_5$
assume either $q=p\equiv\pm1\pmod5$ with $p$ odd prime,
or $q=p^2\equiv-1\pmod5$ with $p$ odd prime. In these cases
$M$ is maximal, has two $G$-classes, and these classes are fused
by $E=\PGL(2,q)$. These are the parameter ranges of the fixed Part II source~\cite{source}.
The classical subgroup facts, as well as the
small subgroup lattices of $A_5$ and $S_4$, are retained premises.

The common two-transitive overgroup lemma is a separate explicit
interface for exhaustivity. It leaves the following possibilities:
\[
\begin{array}{c|c|l}
M&B&\text{possible other parabolic}\\\hline
A_5&D_{10}&D_{20},D_{30},A_5\\
A_5&A_4&S_4,A_5\\
S_4&D_6&D_{12},D_{18},S_4\\
S_4&D_8&D_{16},D_{24},S_4\\
S_4&A_4&A_5
\end{array}
\]
A second $S_4$ over the normal subgroup $A_4$ would have to equal
its unique ambient normalizer and is not a distinct parabolic. The
lists are subject to the actual ambient divisibility conditions. The
common lemma is proved in the separately reviewed manuscript
\emph{Two-transitive coset actions in dimension two}.
All counting arguments below are independent of that exhaustivity step.

\begin{lemma}\label{pair}
Suppose $M$ has two ambient conjugacy classes exchanged by $E$,
contains a single internal class of subgroups of type $B$, and
$\N_M(B)=B$. Assume $M$ is self-normalizing and $B$ is maximal
in $M$. Suppose also that the $B$ subgroups form one $E$-class.
\begin{enumerate}
\item If $B$ has two $G$-classes and $[\N_G(B):B]=2$, then
there are exactly two copies of $M$ above each $B$, both in the
same $G$-class.
\item If $B$ has one $G$-class and $\N_G(B)=B$, then there
are exactly two copies of $M$ above it, one in each $G$-class.
\end{enumerate}
Consequently unordered pairs of distinct copies of $M$ meeting in
$B$ have respectively two or one conjugacy classes, and one class
under $E$ in either case.
\end{lemma}
\begin{proof}
For a fixed ambient class of $M$ containing $B$, its members above
$B$ form a transitive $\N_G(B)$-set, by adjusting a transporter
inside the target $M$. Their stabilizer is $B$. In the first case,
each $M$ class contains only one $B$ class. The outer interchange
of the two classes of each type forces the two $M$ classes to
contain different $B$ classes. Hence the fixed $B$ has exactly two
overgroups in one class and none in the other. In the second case,
each of the two $M$ classes contributes exactly one member.

In both cases there is a unique unordered pair above $B$, and its
intersection is exactly $B$ by maximality of $B$ in $M$. The
intersection recovers $B$ from the pair. Thus conjugacy and $E$
classes of unordered pairs are exactly those of $B$ itself.
Distinct maximal parabolics generate $G$. Their prescribed local
actions are two-transitive, so these pairs give the required geometries.
\end{proof}

\section{Dihedral intersections}
Let $B=D_{2d}$ with $d=3,4$ or $5$. Its characteristic rotation
subgroup lies in a unique torus of $G$, of order
$n=(q-\epsilon)/2$, where $\epsilon=\pm1$ and $d\mid n$.
Write $u=n/d$. In the torus normalizer of $E$ choose
\[
 \langle R,S\mid R^{2n}=S^2=1,\ SRS=R^{-1}\rangle,
 \qquad \N_G(C_d)=\langle R^2,S\rangle,
 \qquad B=\langle R^{2u},S\rangle.
\]
All such $B$ form one $E$-class: the rotations are fixed by the
torus and the possible reflections in $G$ differ by even powers
of $R$, which are conjugate under its full cyclic group.
A rotation $R^k$ normalizes $B$ exactly when $u\mid k$.
Therefore $\N_E(B)=\langle R^u,S\rangle\cong D_{4d}$.
Its intersection with $G$ is all of $D_{4d}$ if $u$ is even,
and is $B$ if $u$ is odd. The $E$-class splits into two $G$-classes
exactly when its normalizer lies in $G$, hence exactly when $u$ is even.

\begin{theorem}\label{dihedral}
Let $(M,B)$ be one of
\[ (A_5,D_{10}),\qquad (S_4,D_6),\qquad(S_4,D_8). \]
The equal-type configurations $(M,M;B)$ have two conjugacy classes
if $u=(q-\epsilon)/(2d)$ is even and one if it is odd. Their
$E$-class count is one. For $k=2$ or $3$, configurations
$(M,D_{2kd};B)$ exist exactly when $kd\mid n$, and then have
two conjugacy classes and one $E$-class.
\end{theorem}
\begin{proof}
Each indicated $B$ is maximal and self-normalizing in $M$, and
forms one internal class there. The equal-type counts follow directly
from Lemma~\ref{pair} and the normalizer calculation above.

A dihedral overgroup $H=D_{2kd}$ must normalize the same rotation
torus: its cyclic rotation subgroup contains $C_d$. If $kd\mid n$,
that torus has a unique subgroup $C_{kd}$, and
$H=\langle C_{kd},B\rangle$ is the unique overgroup of the
required type containing the fixed $B$. Conversely its existence
requires that divisibility. For each ambient class of $M$ there
is one class of chains $(M,B)$, and the unique $H$ over $B$
then gives one configuration. The two classes of $M$ give two
conjugacy classes, interchanged by $E$.

The intersection $M\cap H$ equals $B$: otherwise maximality
of $B$ in $M$ would place $M$ inside a dihedral group, impossible.
Maximality of $M$ then gives generation of $G$. The action of
$D_{4d}$ on cosets of $D_{2d}$ has degree two; that of $D_{6d}$
has degree three with quotient $S_3$. These are two-transitive.
The $M$-side actions are their standard small actions.
\end{proof}

The parity and divisibility tests recover Table 1 rows 5--8 and
12--15, and Table 3 rows 1--8, with the parameter qualifications
below. In Table 3, $6\mid n$, $9\mid n$, $8\mid n$ and
$12\mid n$ are exactly the printed congruences for rows 1,2,5,6.
The full graph groups for the $S_4,S_4,D_8$ and $S_4,D_{16},D_8$
cases are proved in separate accepted manuscripts; these counts do
not substitute for those group-identification arguments.

\section{The intersection $A_4$}
The classical subgroup normalizers give
\[
 \N_G(A_4)=
 \begin{cases}S_4,&q\equiv\pm1\pmod8,\\
 A_4,&q\equiv\pm3\pmod8.
 \end{cases}
\]
In the first case there are two classes of $A_4$, and in the second
there is one; they form one $E$-class. Within $A_5$ the subgroups
$A_4$ form one class and are maximal and self-normalizing.

\begin{theorem}\label{a4}
The configurations $(A_5,A_5;A_4)$ have two conjugacy classes
when $q\equiv\pm1\pmod8$ and one when
$q\equiv\pm3\pmod8$. In both cases they have one $E$-class.
Configurations $(A_5,S_4;A_4)$ exist in the first case only, and
have two conjugacy classes and one $E$-class.
\end{theorem}
\begin{proof}
The equal-type assertion is Lemma~\ref{pair}. For the mixed pair,
a fixed $B=A_4$ has the unique $S_4$ overgroup $N=\N_G(B)$.
The two $A_5$ classes contain different $A_4$ classes, since the
outer coset interchanges each pair of classes. Thus the fixed $B$
belongs to only one class of $A_5$. Its members above $B$ form
a transitive $N$-set with stabilizer $B$, so there are two of them,
interchanged by $N$. Hence there is one mixed configuration for
each of the two $B$ classes, and the outer automorphism interchanges
them. The intersection is $B$; local indices are five and two.
At least the $A_5$ parabolic is maximal, so the graph is connected.
For $q\equiv\pm3\pmod8$, the ambient group contains no $S_4$,
so the mixed configuration does not exist.
\end{proof}

These statements recover Table 1 rows 9--11 and 16--18, and
Table 3 row 9. The mixed $q=9$ case is also Table 6 row 2;
its identical normalizer argument is used in the subfield manuscript.
For these families there is already one class under $E$, so adjoining
field automorphisms cannot change the full automorphism count.

\section{Empty parameter rows and reuse}
In the square-field part of Table 1, the hypotheses impose
$q=p^2\equiv-1\pmod5$ with $p$ odd. Every odd square is one
modulo eight; combining this with $q\equiv4\pmod5$ gives
\[
 q\equiv9\pmod {40}.
\]
Consequently rows 12--14 and 18 have no admissible parameters:
row 12 requires $q\equiv-1\pmod {20}$; row 14 requires
$(q+1)/10$ even; row 18 requires $q\equiv29$ or $19\pmod {40}$.
Row 13 requires $q\equiv-1\pmod {30}$, impossible since a
square is never two modulo three. These are empty conditional
branches, not numerical counterexamples. Rows 15--17 are the actual
square-field configurations. In rows 16--17 the printed alternative
$q\equiv-1\pmod {40}$ is likewise empty, leaving $q\equiv9$.

The remaining cases of Table 1 are already accounted for elsewhere:
rows 2--4 use the accepted scalar-orbit proof in characteristic two;
row 1 uses the unique $D_{30}$ overgroup in the subfield manuscript,
and repeats Table 4 row 5. That row's prime-exponent condition
reduces to $q=16$. The characteristic-five odd-degree case supplies
no geometry, by the subfield uniqueness and exclusion proof. Thus
no A5-family branch is silently omitted.

For source attribution, these are the configurations and numerical
claims of L063 Propositions 5.5 and 5.10, Table 1 and Table 3.
The retained L064 appendix gives the S4 counting argument in
Proposition 20; its proof is recovered by the same formulas. The
source's advertised A5-count appendix pages are not present in that
fixed edition, so the published A5-family proof is joined directly.
The classical subgroup premises and the original constructions remain
attributed. The common normalizer proof and explicit unordered-pair
argument are our derivation. No additional principal refutation,
conjecture answer, Lean verification or compression measurement is
claimed. Whole-original completion still requires the source ledgers
and the other shared proof components.
\begin{thebibliography}{9}
\bibitem{source} F.~Buekenhout, J.~De Saedeleer and D.~Leemans,
\emph{On the rank two geometries of the groups $\PSL(2,q)$: part II},
Ars Mathematica Contemporanea \textbf{6} (2013), 365--388,
Table 1, Table 3 and Propositions 5.5, 5.10.
\bibitem{appendix} The same authors, \emph{Part II (Appendix)},
retained 21-page edition, Proposition 20.
\end{thebibliography}
\end{document}
